\chapter{Security}

\section{The threat model}

Being honest about the threat model keeps the work proportionate. This
application is a single-user site, reachable from the public internet, holding
financial history but no card numbers, no bank credentials and no ability to
move money.

What it must resist:

\begin{itemize}
  \item somebody guessing the password;
  \item somebody reading or forging a session;
  \item a page on another site making the browser act on the user's behalf;
  \item a malicious file or a malicious pasted line;
  \item secrets leaking into the repository.
\end{itemize}

What it does not try to solve: a compromised machine, a hostile hosting
provider, or an attacker who already has the database.

\section{Passwords}

The password is hashed with bcrypt:

\begin{lstlisting}[style=py]
def hash_password(plain: str) -> str:
    return bcrypt.hashpw(plain.encode("utf-8"), bcrypt.gensalt()).decode("utf-8")

def verify_password(plain: str, hashed: str) -> bool:
    try:
        return bcrypt.checkpw(plain.encode("utf-8"), hashed.encode("utf-8"))
    except (ValueError, TypeError):
        return False
\end{lstlisting}

bcrypt is deliberately slow and salts every hash, which is what makes a stolen
database expensive to attack. A general-purpose hash such as SHA-256 is fast,
and speed is exactly the property an attacker wants.

\code{checkpw} compares in constant time, so the verification does not leak
information through timing. The \code{except} catches malformed stored values
and answers "no" rather than raising, which turns a corrupt row into a failed
login instead of a 500 page.

\section{Sessions}

The session is a signed cookie handled by Starlette's session middleware. It
holds one thing: the user id.

\begin{lstlisting}[style=py]
app.add_middleware(
    SessionMiddleware,
    secret_key=settings.secret_key,
    https_only=settings.cookie_secure,
    same_site="lax",
    session_cookie="fm_session",
    max_age=60 * 60 * 12,      # session expires after 12 hours
)
\end{lstlisting}

The cookie is signed, not encrypted, so the contents are readable by the user
but cannot be changed without the key. A user id is not a secret, so that is
fine.

\code{same\_site="lax"} stops the cookie from being attached to cross-site
requests, which removes most of the CSRF surface before the token below is even
considered. \code{https\_only} is tied to configuration so that local
development over plain HTTP still works, while production sets it to true.

And, as chapter \ref{ch:architecture} showed, production refuses to start with
the default signing key. A default key published in a public repository would
let anyone mint a valid session.

\section{CSRF tokens}

Every POST form carries a token that must match one stored in the session:

\begin{lstlisting}[style=py]
def get_csrf_token(request: Request) -> str:
    token = request.session.get(CSRF_SESSION_KEY)
    if not token:
        token = secrets.token_urlsafe(32)
        request.session[CSRF_SESSION_KEY] = token
    return token

def validate_csrf(request: Request, submitted: str | None) -> bool:
    expected = request.session.get(CSRF_SESSION_KEY)
    return bool(expected) and bool(submitted) and secrets.compare_digest(expected, submitted)
\end{lstlisting}

The token is generated on first use and reused for the session, which is simpler
than rotating it per form and enough given the same-site cookie policy.

Handlers that fail validation redirect instead of raising. A stale tab that
submits an old token sends the user back to the page rather than to an error.

\section{Slowing down brute force}

The login form is the one endpoint an attacker can hammer. A sliding window per
IP address limits it:

\begin{lstlisting}[style=py]
class LoginRateLimiter:
    def __init__(self, max_attempts: int = 8, window_seconds: int = 900):
        ...
    def _prune(self, ip: str, now: float) -> list[float]:
        recent = [t for t in self._fails.get(ip, []) if now - t < self.window]
        self._fails[ip] = recent
        return recent
\end{lstlisting}

Eight failures in fifteen minutes and the form answers \code{429}. Only failures
are recorded, so a legitimate user is never affected. The state is a dictionary
protected by a lock, which is appropriate for a single-process application; a
restart clears it, and that is an acceptable trade for having no extra service
to run.

Behind a proxy the socket address is the proxy, so the client address is read
from \code{X-Forwarded-For} when present.

\section{Response headers}

One middleware sets the same headers on every response:

\begin{lstlisting}[style=py]
response.headers["X-Content-Type-Options"] = "nosniff"
response.headers["X-Frame-Options"] = "DENY"
response.headers["Referrer-Policy"] = "no-referrer"
response.headers["Permissions-Policy"] = "geolocation=(), microphone=(), camera=()"
response.headers["Content-Security-Policy"] = (
    "default-src 'self'; "
    "script-src 'self' 'unsafe-inline' https://cdn.jsdelivr.net https://unpkg.com; "
    ...
    "connect-src 'self'; frame-ancestors 'none'; base-uri 'self'")
if settings.cookie_secure:
    response.headers["Strict-Transport-Security"] = "max-age=31536000; includeSubDomains"
\end{lstlisting}

\code{X-Frame-Options: DENY} prevents clickjacking, \code{nosniff} stops content
type guessing, and the CSP restricts where scripts and styles may come from. The
policy still allows \code{'unsafe-inline'} for scripts, because the chart
initialisation lives in the templates; that is the weakest line in the policy and
the honest place to improve it later.

HSTS is only sent when the deployment is on HTTPS, since sending it from a plain
HTTP development server would pin the browser to a scheme the server does not
speak.

\section{Uploads and pasted text}

Statement uploads are limited in three ways: an eight megabyte cap, an accepted
extension per broker, and a parser that never raises. The size check happens
after reading, which is acceptable for a single-user application where the only
uploader is the owner.

The pasted-arithmetic parser described in chapter 4 is a security control as
much as a convenience: it is the reason a line of text from a phone cannot run
code.

\section{Secrets and the repository}

The last control is about the repository rather than the code. \code{.env} is
ignored by git and only \code{.env.example} is committed, with placeholder
values. Nothing in the source contains a name, an email, an account or a path
belonging to a real person; every one of those comes from the environment.

The API documentation is hidden when \code{IS\_PROD=true}, which removes a free
map of the application from anyone who finds the URL.

\begin{keypoint}[The pattern worth copying]
Each control here is a few lines. What makes them effective is that they are
applied in one place: one middleware for headers, one function for CSRF, one
settings object for secrets. Security spread across handlers is security that
will be forgotten in the next handler.
\end{keypoint}
